% =====================================================================
% 1. ENCODAGE, POLICES ET LANGUE
% =====================================================================
\usepackage[utf8]{inputenc}       % Gestion de l'UTF-8
\usepackage[T1]{fontenc}          % Encodage des caractères pour le français
\usepackage[french]{babel}        % Support du français
\usepackage{lmodern}              % Police propre et vectorielle
\usepackage{gensymb}

% =====================================================================
% 2. MISE EN PAGE ET GÉOMÉTRIE
% =====================================================================
\usepackage[a4paper, margin=3.2cm, top=3.5cm, bottom=4cm]{geometry}
\usepackage{setspace}             % Gestion de l'interligne
\usepackage{fancyhdr}             % En-têtes et pieds de page personnalisés (Remplace fancyheadings)
\usepackage{lastpage}             % Pour obtenir le numéro de la dernière page (\pageref{LastPage})
\usepackage[absolute,overlay]{textpos} % Positionnement absolu de blocs de texte
\usepackage{eso-pic}              % Pour ajouter des éléments en arrière-plan de la page

% =====================================================================
% 3. MATHÉMATIQUES ET SCIENCES
% =====================================================================
\usepackage{amsmath}              % Commandes mathématiques de base
\usepackage{amsfonts}             % Polices mathématiques
\usepackage{amssymb}              % Symboles mathématiques
\usepackage{siunitx}              % Alignement et écriture des unités scientifiques

% =====================================================================
% 4. GRAPHIQUES, TRACÉS ET DIAGRAMMES
% =====================================================================
\usepackage{graphicx}             % Inclusion d'images
\graphicspath{{Param/}, {Images/}} % Dossiers où chercher les images

\usepackage{tikz}                 % Dessin vectoriel personnalisé
\usetikzlibrary{shapes.geometric, arrows, chains, positioning, babel}
\usepackage{pgfplots}             % Tracé de graphiques mathématiques de haute qualité
\pgfplotsset{compat=1.18}

\usepackage[european, straightvoltages, RPvoltages]{circuitikz} % Tracé de circuits électriques

% =====================================================================
% 5. TABLEAUX
% =====================================================================
\usepackage{booktabs}             % Tableaux de qualité professionnelle (\toprule, \midrule, etc.)
\usepackage{tabularx}             % Tableaux à largeur fixe adaptative
\usepackage{array}                % Personnalisation avancée des colonnes
\usepackage{longtable}            % Tableaux multipages
\usepackage{makecell}             % Retours à la ligne et mise en forme dans les cellules
\usepackage{rotating}             % Pour tourner les tableaux ou les figures (mode paysage local)

% =====================================================================
% 6. STRUCTURATION, BLOCS ET COULEURS
% =====================================================================
\usepackage[table]{xcolor}        % Gestion des couleurs et coloration des lignes de tableaux
\usepackage{titlesec}             % Personnalisation des titres (\section, \subsection)
\usepackage{enumitem}             % Personnalisation des listes à puces
\usepackage{subcaption}           % Légendes pour les sous-figures
\usepackage{float}                % Forcer le positionnement des éléments graphiques ([H])
\usepackage{listings}             % Inclusion et coloration de code source (Python, C++, etc.)
\usepackage[most]{tcolorbox}      % Boîtes d'alerte et encadrés stylisés
\usepackage{pifont}               % Symboles et puces spécifiques (\ding)
\usepackage{lipsum}               % Générateur de faux texte de test

% =====================================================================
% 7. LIENS HYPERTEXTES (Toujours charger en dernier !)
% =====================================================================
\usepackage[hidelinks]{hyperref}  % Liens cliquables internes et externes sans bordures rouges